\documentclass[11pt]{article}
\usepackage[margin=1in]{geometry}
\usepackage{amsmath,amssymb,amsthm}
\newtheorem{theorem}{Theorem}
\newcommand{\Id}{\mathbb I}
\begin{document}
\begin{center}\large An exact order-seven cone for a two-parameter C2 surface\end{center}

Use the canonical slack
\[
 f_c=d\bigl[1-(E_{00}-E_{01}+E_{10}+E_{11})/2\bigr]
           +4\sum_i a_i p_i,
\]
where $c=(d,\lambda,n,m,C,A,D,\alpha,\beta)$ and the eight events are
\[
 (++|01),(+-|00),(+-|11),(-+|10),(-+|00),(-+|11),(--|01),(+-|10).
\]
The nonnegative source coordinates and the requirement $f_c\geq0$ on all
quantum behaviours define the canonical valid cone.

\paragraph{The stationary sources and their exact domain.}
Set $d=1$, $\lambda=\beta=0$, and let
\[
 0<r<1,\qquad 0<z\leq\frac{1-r}{r(1+3r)},\qquad
 T=1+(1+2r)z.
\]
The six contact equations at
\[
 u=v=\sqrt z,\qquad
 \psi=(1,-r\sqrt z,r\sqrt z,1)^T
\]
give
\begin{align*}
 n&=\frac{1-r-r(1+3r)z}{4rT},&
 m&=\frac{(1-r)(1+rz)}{4rT},&
 C&=\frac{r(1+rz)}{(1-r)T},\\
 A&=\frac{(1+r)(1+rz)}{4rT},&
 D&=\frac{(1+r)(1-rz)}{4rT},&
 \alpha&=\frac{r(1+r)z^2}{(1-rz)T}.
                                                               \tag{1}
\end{align*}
Here $n\geq0$ is equivalent to the stated upper bound; that bound implies
$rz\leq(1-r)/(1+3r)<1$, so all other required event weights are positive.
The six equations have a nonzero minor, deleting the $d$ column,
\[
 -\frac{32768c^6r^4(r-1)(c^2r-1)(2c^2r+c^2+1)}{(1+c^2)^5},
 \qquad c=\sqrt z.
\]
The matrix here has columns $(d,n,m,C,A,D,\alpha)$ and rows
$(\mathcal B\psi)_1,\ldots,(\mathcal B\psi)_4,
\psi^T\mathcal B_u\psi,\psi^T\mathcal B_v\psi$, evaluated at $u=v=c$,
where $\mathcal B$ is the uncleared Bell operator and the unnormalised
state $\psi$ is held fixed. The first four rows are not doubled gradients.
Thus (1) is the unique source with $d=1$ for this contact.
Once validity is established below, it also proves canonical extremality:
in any sum of valid nonnegative canonical sources, $\lambda=\beta=0$
and the contact value persist in every summand. Positivity then forces
the four kernel equations, while the interior angular minimum forces
the two derivative equations. The one-dimensional solution space makes
every summand a nonnegative multiple of (1), including at $n=0$.

Put $x=rz$, $\kappa=r+x+2rx$. The entire domain becomes
\[
 r,x>0,\qquad r+x+3rx\leq1,                              \tag{2}
\]
which implies $r,x<1$. Multiplication by $4\kappa(1-r)(1-x)>0$
gives the polynomial source, written only in its seven active coordinates,
\[
 s(r,x)=\begin{pmatrix}
 4\kappa(1-r)(1-x)\\
 (1-r-x-3rx)(1-r)(1-x)\\
 (1-r)^2(1-x^2)\\
 4r^2(1-x^2)\\
 (1-r^2)(1-x^2)\\
 (1-r^2)(1-x)^2\\
 4x^2(1-r^2)
 \end{pmatrix}_{(d,n,m,C,A,D,\alpha)}.                   \tag{3}
\]
Every coordinate has degree at most two in each parameter.

\begin{theorem}
Let $K$ be the closed conic hull of (3) over (2). It is exactly the cone
specified by $\lambda=\beta=0$, the linear identities
\[
 d=4A-2m-2D,\qquad n=m+D-A,                              \tag{4}
\]
and the fixed linear matrix inequalities
\[
 \begin{pmatrix}m&A-m\\A-m&C\end{pmatrix}\succeq0,
 \qquad
 \begin{pmatrix}D&A-D\\A-D&\alpha\end{pmatrix}\succeq0,
                                                               \tag{5}
\]
\[
 A-m\geq0,\qquad A-D\geq0,\qquad m+D-A\geq0.            \tag{6}
\]
This has total PSD order seven. Every source in $K$ is quantum-valid.
\end{theorem}

\begin{proof}[Proof of the exact source cone]
On the surface, (4) holds and both matrices in (5) have rank one:
\[
 (A-m)^2=mC,\qquad (A-D)^2=D\alpha.
\]
The inequalities (6) follow immediately from (3) and (2). The cone
defined by (4)--(6) is closed, so it contains $K$.

Conversely, suppose (4)--(6) hold. If either $m=0$ or $D=0$, its
positive semidefinite block forces $A=0$, and (6) forces $m=D=0$.
Only $C,\alpha\geq0$ remain. Both corresponding event rays lie in $K$:
\[
 \lim_{\varepsilon\downarrow0}s(1-\varepsilon,\varepsilon/8)
       =4e_C,\qquad
 \lim_{\varepsilon\downarrow0}s(\varepsilon/8,1-\varepsilon)
       =4e_\alpha.
                                                               \tag{7}
\]
These paths satisfy (2).

Otherwise $A,m,D>0$. The two Schur complements give
\[
 C\geq C_0=\frac{(A-m)^2}{m},\qquad
 \alpha\geq\alpha_0=\frac{(A-D)^2}{D}.
\]
Set
\[
 r=\frac{A-m}{A+m},\qquad x=\frac{A-D}{A+D}.
\]
Then $0\leq r,x<1$, and
\[
 r+x+3rx-1=\frac{4A(A-m-D)}{(A+m)(A+D)}\leq0.
\]
The source with $C=C_0$, $\alpha=\alpha_0$ and all other coordinates
unchanged is precisely
\[
 \frac{A}{(1-r^2)(1-x^2)}s(r,x).
\]
Boundary points $r=0$ or $x=0$, including $r=x=0$, are limits of (2).
In particular, the latter gives $d=0$, $n=m=A=D$, $C_0=\alpha_0=0$.
Adding $(C-C_0)e_C+(\alpha-\alpha_0)e_\alpha$ and using (7) puts the
original source in $K$. This proves both inclusions with all boundary
cases explicit. It also shows that (4)--(6) imply all canonical source
signs: the two PSD blocks give $m,C,D,\alpha\geq0$, (6) gives $A,n\geq0$,
and $m,D\leq A$ gives $d\geq0$.
\end{proof}

\paragraph{A standalone physical-validity proof.}
For the scaled source (3), choose
\[
 A_0=B_0=\sigma_z,\qquad
 A_1=\frac{(1-u^2)\sigma_z+2u\sigma_x}{1+u^2},\qquad
 B_1=\frac{(1-v^2)\sigma_z+2v\sigma_x}{1+v^2}.
\]
Let $P(u,v)$ be its Bell-operator state matrix multiplied by
$D_uD_v=(1+u^2)(1+v^2)$. For $uv\ne0$ its $Z_s=0$ state block is
positive definite: its terms proportional to $d$ alone equal
\[
 2d\bigl[(uW+uvX_s-Y)^2+D_uv^2Y^2+D_vX_s^2\bigr].       \tag{8}
\]
All event contributions are positive semidefinite. Thus $P$ has at least
three positive eigenvalues, and determinant nonnegativity suffices for
positivity of the whole matrix.

Write $z=x/r$, $\mathsf X=u^2/z$, $\mathsf Y=v^2/z$ and
$G=16(1-r^2)^3(1-x^2)^3$. Direct expansion of this four-by-four state
matrix gives
\[
 \det P=16G D_u^2D_v^2u^2v^2\mathcal R(\mathsf X,\mathsf Y),
                                                               \tag{9}
\]
where the entire Laurent factor is
\begin{align*}
\mathcal R={}&a(\mathsf X+\mathsf X^{-1}-2)
 +b(\mathsf Y+\mathsf Y^{-1}-2)\\
 &+c_0(\mathsf X/\mathsf Y+\mathsf Y/\mathsf X-2)
 +\Delta(\mathsf Y/\mathsf X+\mathsf X+\mathsf Y^{-1}-3),\\
a={}&8rx(1-r),\qquad b=8rx(1-x),\\
c_0={}&(1+r)(1+x)(1-r-x-3rx),\qquad \Delta=8rx(r+x).
                                                               \tag{10}
\end{align*}
All four coefficients are nonnegative on (2). The first three brackets
are reciprocal-pair AM--GM inequalities. The final bracket is three-term
AM--GM, since the product of its three terms is one. Thus (8)--(10) prove
positivity for $uv\ne0$. Continuity of the uncleared Bell matrix includes
zero and infinite half-angle coordinates. Two-reflection block
decomposition then proves validity in arbitrary dimension; the real state
matrix being positive also covers complex states. Closure and conic
combination preserve validity, so the whole cone in (4)--(6) is sound.

The accompanying exact script reconstructs all six contact equations,
the source scaling, both rank-one identities, the inverse decomposition,
the two event-ray limits, and the full four-variable polynomial identity
(9). No varying testing-function family or parameter grid enters the cone.
\end{document}
